%% ============================================================
%% ENTREGABLE — TE3001B Fundamentación de Robótica
%% Teleoperación con Sensado de Fuerza: xArm Lite 6 (Maestro-Esclavo)
%% Formato IEEE Transactions — Plantilla para Alumnos
%% Prof. Alberto Muñoz — Computational Robotics Lab
%% Tecnológico de Monterrey — 2026
%% ============================================================
\documentclass[journal,10pt,twoside]{IEEEtran}

%% ── Paquetes base ────────────────────────────────────────────
\usepackage[utf8]{inputenc}
\usepackage[T1]{fontenc}
\usepackage[spanish,es-nodecimaldot]{babel}
%% Restore IEEE English labels (babel[spanish] overrides them)
\renewcommand{\abstractname}{Abstract}
\def\IEEEkeywordsname{Index Terms}
\usepackage{amsmath,amssymb,bm}
\usepackage{graphicx}
\usepackage{xcolor}
\usepackage{booktabs}
\usepackage{multirow}
\usepackage{array}
\usepackage{colortbl}
\usepackage{caption}
\usepackage{subcaption}
\usepackage{listings}
\usepackage{tcolorbox}
\usepackage{enumitem}
\usepackage{tikz}
\usepackage{pgfplots}
\usepackage{url}
\usepackage{hyperref}   %% hyperref DESPUÉS de xcolor para evitar conflictos
\pgfplotsset{compat=1.18}
\usetikzlibrary{arrows.meta,shapes,positioning,calc,backgrounds}
\tcbuselibrary{breakable,skins}

%% ── Colores institucionales ──────────────────────────────────
%% IMPORTANTE: definir colores ANTES de \hypersetup
\definecolor{tec_blue}{RGB}{0,68,124}
\definecolor{tec_lightblue}{RGB}{0,140,200}
\definecolor{tec_red}{RGB}{190,30,45}
\definecolor{tec_gray}{RGB}{100,100,100}
\definecolor{code_bg}{RGB}{245,247,250}
\definecolor{code_kw}{RGB}{0,100,180}
\definecolor{code_str}{RGB}{180,50,20}
\definecolor{code_cmt}{RGB}{80,140,80}
\definecolor{box_todo}{RGB}{255,243,205}
\definecolor{box_todo_border}{RGB}{200,140,0}
\definecolor{box_rubric}{RGB}{232,245,233}
\definecolor{box_rubric_border}{RGB}{67,160,71}
\definecolor{box_theory}{RGB}{227,242,253}
\definecolor{box_theory_border}{RGB}{25,118,210}

%% hypersetup aquí, DESPUÉS de definir tec_blue ───────────────
\hypersetup{colorlinks=true, linkcolor=tec_blue, urlcolor=tec_blue,
            citecolor=tec_blue}

%% ── Cajas de tcolorbox ───────────────────────────────────────
\newtcolorbox{todobox}[1][Completar]{
  colback=box_todo, colframe=box_todo_border,
  title={\textbf{$\triangleright$ #1}},
  fonttitle=\bfseries\color{tec_gray},
  breakable, enhanced, left=4pt, right=4pt
}
\newtcolorbox{rubricbox}[1][Rúbrica]{
  colback=box_rubric, colframe=box_rubric_border,
  title={\textbf{\checkmark\ Criterio de Evaluación — #1}},
  fonttitle=\bfseries\small\color{box_rubric_border},
  breakable, enhanced, left=4pt, right=4pt
}
\newtcolorbox{theorybox}[1][Marco Teórico]{
  colback=box_theory, colframe=box_theory_border,
  title={\textbf{(i) Contexto: #1}},
  fonttitle=\bfseries\small\color{tec_blue},
  breakable, enhanced, left=4pt, right=4pt
}

%% ── Estilo de código ─────────────────────────────────────────
\lstdefinestyle{python}{
  language=Python,
  backgroundcolor=\color{code_bg},
  basicstyle=\ttfamily\scriptsize,
  keywordstyle=\color{code_kw}\bfseries,
  stringstyle=\color{code_str},
  commentstyle=\color{code_cmt}\itshape,
  numbers=left, numberstyle=\tiny\color{tec_gray},
  numbersep=6pt, frame=single, tabsize=4,
  showstringspaces=false, breaklines=true,
  morekeywords={import,from,as,def,class,return,True,False,None,self},
  literate={á}{{\'{a}}}1 {é}{{\'{e}}}1 {í}{{\'{i}}}1
           {ó}{{\'{o}}}1 {ú}{{\'{u}}}1 {ñ}{{\~{n}}}1,
}
\lstset{style=python}

%% ── Macros de ayuda ──────────────────────────────────────────
\newcommand{\answerline}[1][3cm]{\underline{\hspace{#1}}}
\newcommand{\pts}[1]{\textcolor{tec_red}{\textbf{[#1 pts]}}}
\newcommand{\TODO}[1]{\begin{todobox}[Respuesta del equipo]
  \textit{#1}
\end{todobox}}
\newcommand{\FIGURE}[1]{\begin{todobox}[Insertar figura]
  \centering
  %% \resizebox escala el tikzpicture al ancho disponible
  \resizebox{\linewidth}{!}{%
    \begin{tikzpicture}
      \draw[tec_gray, dashed, line width=1pt] (0,0) rectangle (8,4);
      \node[tec_gray, font=\small] at (4,2) {#1};
    \end{tikzpicture}%
  }
\end{todobox}}

%% ── Metadatos del artículo ───────────────────────────────────
\begin{document}

\title{Teleoperación con Sensado de Fuerza Física\\
  entre dos xArm Lite 6:\\
  Diseño, Implementación y Evaluación Experimental}

\author{%
  \IEEEauthorblockN{%
    Gerardo Fregoso Jiménez\IEEEauthorrefmark{1},\;
    Luis Fernando Ruíz Neria\IEEEauthorrefmark{1},\;
    Daniel Eduardo Hinojosa Alvarado\IEEEauthorrefmark{1},\;
    José Raúl Arredondo López\IEEEauthorrefmark{1}%
  }\\
  \IEEEauthorblockA{%
    \IEEEauthorrefmark{1}%
    Computational Robotics Lab, Tecnológico de Monterrey,
    Monterrey, N.L., México\\
    \texttt{\{A00838645, A00839528, A00838156, A01742153\}@tec.mx}
  }
}

\maketitle

%% ── Abstract ─────────────────────────────────────────────────
\begin{abstract}
This paper presents the design, implementation, and experimental
evaluation of a bilateral teleoperation system with force sensing
between two xArm~Lite~6 six-degree-of-freedom manipulators. The
master-slave architecture operates at 100~Hz over UDP communication
using binary packets of 56 and 104~bytes, achieving an effective
delay of $T_d < 2$~ms on a local area network. Cartesian force at
the slave end-effector is estimated without an external F/T sensor
by computing joint-torque residuals: the damped least-squares
pseudo-inverse of the geometric Jacobian ($\lambda = 0.01$) maps
gravity-compensated residual torques to a six-component wrench,
followed by a first-order low-pass filter at $f_c = 5$~Hz. Since
the xArm Python SDK does not support direct joint-torque injection
in user mode, haptic feedback to the master is implemented as a
joint-space position-impedance correction with force-reflection
gain $\beta = 0.3$. An additional physical force sensor based on an
ESP32 microcontroller with micro-ROS was designed and built,
measuring three-axis forces via an analog joystick ($\pm3$~N) and a
spring-ultrasonic assembly ($k = 1000$~N/m), publishing
\texttt{WrenchStamped} messages to ROS~2 at 50~Hz over WiFi.
Experiments demonstrate functional teleoperation with joint tracking
RMS error $<$~[X]° and force estimation error $<$~[X]~N against a
precision scale reference.
\end{abstract}

\begin{IEEEkeywords}
  teleoperación, sensado de fuerza, xArm Lite 6, control de impedancia,
  retroalimentación háptica, micro-ROS, admitancia
\end{IEEEkeywords}

%% ============================================================
\section{Introducción}
%% ============================================================

\IEEEPARstart{L}{a} teleoperación con retroalimentación de fuerza
constituye un paradigma fundamental en robótica cuando se requiere que
un operador humano interactúe con entornos remotos o peligrosos sin
presencia física~\cite{murray1994}. El xArm Lite 6, desarrollado por
UFACTORY~\cite{xarm_github}, es un manipulador de 6 grados de libertad
(DOF) con capacidad de control de torque que lo hace idóneo para experimentos
de sensado de fuerza por torque articular.

\begin{theorybox}[Motivación]
  El reto central de este proyecto es cerrar el lazo de fuerza entre
  dos brazos robóticos: el robot \textbf{Maestro}, operado por el humano,
  y el robot \textbf{Esclavo}, que interactúa con el entorno.
  La fuerza de contacto estimada en el esclavo debe reflejarse como
  torque en las articulaciones del maestro, creando la sensación
  háptica que guía al operador.
\end{theorybox}

\subsection{Contexto y Motivación}

La teleoperación con sensado de fuerza permite extender la capacidad
sensorial del operador humano más allá de sus limitaciones físicas, ya
sea en distancia, escala o condiciones ambientales adversas. En el
contexto industrial, la manipulación remota de materiales peligrosos
(residuos nucleares, explosivos) requiere que el operador perciba las
fuerzas de contacto en tiempo real para evitar daños al entorno o al
propio manipulador~\cite{niemeyer2008}. Por ejemplo, en plantas de
desmantelamiento nuclear, sistemas bilaterales equipados con sensores
F/T de muñeca operan con retardos de red de decenas de milisegundos,
imponiendo severas restricciones de estabilidad que son abordadas
mediante control de onda~\cite{niemeyer2008}.

En el ámbito médico, la cirugía robótica mínimamente invasiva (e.g.,
sistemas tipo da Vinci) depende de la retroalimentación háptica para
que el cirujano calibre la fuerza de tracción sobre tejidos blandos.
Estudios clínicos han demostrado que la ausencia de retroalimentación
de fuerza incrementa significativamente la tasa de errores en tareas de
sutura~\cite{siciliano2009}. La estimación de fuerza a partir de
torques articulares —sin sensor externo— abre la posibilidad de dotar
de haptics a sistemas quirúrgicos de bajo costo.

Desde el punto de vista académico, el problema de teleoperación con
retardo es un banco de pruebas canónico para teorías de estabilidad de
sistemas interconectados~\cite{colgate1993,spong2006}. La
implementación con plataformas accesibles como el xArm Lite 6 y el
ecosistema ROS~2 permite replicar y extender resultados del estado del
arte de forma reproducible.

\subsection{Contribuciones del Trabajo}

Las contribuciones específicas de esta implementación son:
\begin{itemize}[itemsep=0.2em]
  \item \textbf{Protocolo UDP de baja latencia:} diseño de un protocolo
        binario de comunicación a 100~Hz con paquetes de 56 y 104~bytes
        y medición de RTT, jitter y percentil 99 de latencia.
  \item \textbf{Estimador de fuerza sensorless:} implementación del
        estimador $\hat{\bm{F}}_{ext} = (\bm{J}^\top)^+
        (\bm{\tau}_{meas} - \bm{\tau}_{grav})$ con pseudoinversa
        amortiguada (DLS) y filtro pasa-bajas a 5~Hz, incluyendo
        calibración automática del nivel de ruido en reposo.
  \item \textbf{Retroalimentación háptica por impedancia de posición:}
        alternativa al control de torque directo que no está disponible
        en el SDK Python del xArm, implementando la corrección
        $\delta\bm{q} = -k_h \bm{J}^\top \beta \hat{\bm{F}}_{ext}$.
  \item \textbf{Sensor de fuerza físico con micro-ROS:} diseño y
        construcción de un sensor de tres ejes basado en ESP32, joystick
        analógico y resorte-ultrasónico, integrado al sistema ROS~2
        mediante micro-ROS vía WiFi a 50~Hz.
  \item \textbf{Integración ROS~2 + MoveIt Servo:} nodos de admitancia
        (maestro) e impedancia (esclavo) que publican
        \texttt{TwistStamped} a MoveIt~Servo, habilitando teleoperación
        en espacio cartesiano a 50~Hz.
\end{itemize}

\subsection{Estructura del Artículo}
El resto de este artículo se organiza como sigue:
la Sección~\ref{sec:background} presenta los fundamentos teóricos;
la Sección~\ref{sec:architecture} describe la arquitectura del sistema;
la Sección~\ref{sec:force} detalla la estrategia de sensado de fuerza;
la Sección~\ref{sec:control} expone el esquema de control;
la Sección~\ref{sec:implementation} documenta la implementación con xArm SDK;
la Sección~\ref{sec:experiments} reporta los resultados experimentales;
la Sección~\ref{sec:discussion} discute los hallazgos;
y la Sección~\ref{sec:conclusions} concluye el trabajo.

%% ============================================================
\section{Marco Teórico}
\label{sec:background}
%% ============================================================

\subsection{Cinemática y Jacobiano del xArm Lite 6}

El xArm Lite 6 es un manipulador de $n=6$ DOF con articulaciones
rotacionales. Su cinemática directa se describe mediante la convención
Denavit-Hartenberg (DH) o Product of Exponentials (PoE).

\begin{theorybox}[Jacobiano Geométrico]
  El Jacobiano $\bm{J}(\bm{q}) \in \mathbb{R}^{6\times6}$ mapea
  velocidades articulares a velocidades cartesianas del efector final:
  \[
    \begin{pmatrix} \bm{v} \\ \bm{\omega} \end{pmatrix} = \bm{J}(\bm{q})\,\dot{\bm{q}}
  \]
  Su transpuesta mapea fuerzas/torques cartesianos a torques articulares:
  \[
    \bm{\tau} = \bm{J}^\top(\bm{q})\,\bm{F}_{ext}
  \]
  Esta relación es la base del \emph{torque-based force sensing}.
\end{theorybox}

\subsubsection*{Pregunta~T1 \pts{8}}
\begin{enumerate}[label=(\alph*)]
  \item Expliquen la diferencia entre el \textbf{Jacobiano analítico}
        y el \textbf{Jacobiano geométrico} para el xArm Lite 6.
        ¿Cuál utilizan en su implementación y por qué?
  \item Identifiquen y describan \textbf{dos configuraciones singulares}
        del xArm Lite 6 a partir de su geometría DH. ¿Qué ocurre con la
        estimación de fuerza $\hat{\bm{F}} = (\bm{J}^\top)^+\bm{\tau}$
        cerca de esas singularidades?
  \item ¿Cómo detectan programáticamente que el robot está cerca de
        una singularidad? Escriban la condición en términos del
        número de condición $\kappa(\bm{J})$.
\end{enumerate}

\textbf{(a) Jacobiano analítico vs.\ geométrico.}
El \textbf{Jacobiano analítico} $\bm{J}_a \in \mathbb{R}^{6\times6}$
se define como la derivada parcial de la pose del efector final
$\bm{x} = [\bm{p}^\top, \bm{\phi}^\top]^\top$ (donde $\bm{\phi}$ es
una representación mínima de la orientación, e.g., ángulos de Euler
RPY) con respecto a las coordenadas articulares:
$\dot{\bm{x}} = \bm{J}_a(\bm{q})\dot{\bm{q}}$.

El \textbf{Jacobiano geométrico} $\bm{J} \in \mathbb{R}^{6\times6}$
relaciona directamente $\dot{\bm{q}}$ con la velocidad lineal
$\bm{v}$ y angular $\bm{\omega}$ del efector en el frame base:
\begin{equation}
  \bm{J}(\bm{q}) =
  \begin{pmatrix} \bm{J}_v \\ \bm{J}_w \end{pmatrix},\;
  \text{col.\ } i{:}\;
  \begin{cases}
    \bm{J}_{v,i} = \bm{z}_{i-1} \!\times\! (\bm{p}_e - \bm{p}_{i-1}) \\
    \bm{J}_{w,i} = \bm{z}_{i-1}
  \end{cases}
  \label{eq:geom_jac}
\end{equation}
donde $\bm{z}_{i-1}$ y $\bm{p}_{i-1}$ son el eje~$z$ y el origen del
frame $(i{-}1)$ en el frame base.

En nuestra implementación se utiliza el \textbf{Jacobiano geométrico}
(función \texttt{geometric\_jacobian(q)} en \texttt{kinematics.py}).
La razón es doble: (i) es directamente aplicable en la relación de
estática $\bm{\tau} = \bm{J}^\top \bm{F}$ sin necesidad de ningún
factor de transformación adicional; (ii) evita las singularidades de
representación de la orientación que presentan los ángulos de Euler
(e.g., gimbal lock para RPY cuando $\phi_2 = \pm\pi/2$).

Los parámetros DH del xArm Lite 6, extraídos del manual UFACTORY e
implementados en \texttt{jacobian\_utils.py}, se muestran en la
Tabla~\ref{tab:dh}.

\begin{table}[h]
  \centering
  \caption{Parámetros DH estándar del xArm Lite 6}
  \label{tab:dh}
  \small
  \begin{tabular}{ccccc}
    \toprule
    \textbf{Joint} & $a_i$ (m) & $d_i$ (m) & $\alpha_i$ (rad) & $\theta_{off}$ (rad) \\
    \midrule
    1 & 0.0000 & 0.2433 & $+\pi/2$ & $0$       \\
    2 & 0.2000 & 0.0000 & $0$      & $-\pi/2$  \\
    3 & 0.0870 & 0.0000 & $+\pi/2$ & $0$       \\
    4 & 0.0000 & 0.2276 & $-\pi/2$ & $0$       \\
    5 & 0.0000 & 0.0000 & $+\pi/2$ & $0$       \\
    6 & 0.0000 & 0.0615 & $0$      & $0$       \\
    \bottomrule
  \end{tabular}
\end{table}

\textbf{(b) Configuraciones singulares.}
Para el xArm Lite 6 se identifican dos configuraciones singulares
principales:
\begin{enumerate}[label=\roman*.]
  \item \textbf{Singularidad de codo (\textit{elbow}):} cuando
        $q_3 \approx 0°$, los eslabones~2 y~3 quedan colineales (el
        brazo se extiende completamente). El Jacobiano pierde rango en
        la dirección radial: la articulación~3 ya no aporta movimiento
        independiente al efector, y $\det(\bm{J}) \to 0$. Cerca de esta
        configuración, la pseudoinversa $(\bm{J}^\top)^+$ amplifica
        enormemente el ruido de torque, produciendo estimaciones
        $\hat{\bm{F}}_{ext}$ espurias de cientos de N.
  \item \textbf{Singularidad de muñeca (\textit{wrist}):} cuando
        $q_5 \approx 0°$, los ejes de las articulaciones~4 y~6 se
        alinean. El robot pierde un grado de libertad rotacional:
        sólo puede girar alrededor de la dirección común de esos dos
        ejes. La estimación de torques de momento ($M_x, M_y$) se
        vuelve indefinida.
\end{enumerate}

\FIGURE{Fig. T1: Configuraciones singulares del xArm Lite 6.
Señalen las articulaciones involucradas.}

\textbf{(c) Detección de singularidad.}
La singularidad se detecta mediante el número de condición del
Jacobiano:
\begin{equation}
  \kappa(\bm{J}) = \frac{\sigma_{max}}{\sigma_{min}}
  \label{eq:kappa}
\end{equation}
donde $\sigma_{max}$ y $\sigma_{min}$ son el mayor y menor valor
singular de $\bm{J}$. Implementado en \texttt{condition\_number(J)}
como \texttt{np.linalg.cond(J)}. En nuestro sistema se utilizan dos
umbrales:
\begin{align*}
  \kappa &> 50  \;\Rightarrow\; \text{suspender háptica en maestro} \\
  \kappa &> 200 \;\Rightarrow\; \text{zona singular severa}
\end{align*}
El umbral $\kappa_{th} = 50$ fue seleccionado empíricamente para
garantizar que $\hat{\bm{F}}_{ext}$ tenga error $< 5$~N en las
regiones de trabajo habituales.

\subsection{Dinámica del Manipulador}

La dinámica de un robot rígido de $n$ DOF se expresa como:
\begin{equation}
  \bm{M}(\bm{q})\ddot{\bm{q}} + \bm{C}(\bm{q},\dot{\bm{q}})\dot{\bm{q}}
  + \bm{g}(\bm{q}) = \bm{\tau} - \bm{J}^\top\bm{F}_{ext}
  \label{eq:dynamics}
\end{equation}
donde $\bm{M}$ es la matriz de inercia, $\bm{C}$ la matriz de
Coriolis/centrífuga y $\bm{g}$ el vector gravitacional.

\subsubsection*{Pregunta~T2 \pts{10}}
\begin{enumerate}[label=(\alph*)]
  \item Partiendo de la Ec.~\eqref{eq:dynamics}, \textbf{deriven} la
        ley de control de par calculado (Computed Torque Control, CTC)
        y demuestren que produce la dinámica de error lineal:
        \[
          \ddot{\bm{e}} + K_v\dot{\bm{e}} + K_p\bm{e} = \bm{0}
        \]
        ¿Qué condiciones sobre $K_p$ y $K_v$ garantizan
        estabilidad asintótica global?
  \item Para el xArm Lite 6 operando en modo \texttt{SERVO} de
        la SDK, el bucle de control interno trabaja a 250~Hz.
        Expliquen cómo esto limita la elección de $K_p$ y $K_v$.
  \item Expresen los polos del sistema cerrado en términos de
        $\omega_n$ (frecuencia natural) y $\zeta$ (amortiguamiento).
        ¿Qué valor de $\zeta$ usaron y por qué?
\end{enumerate}

\textbf{(a) Derivación del CTC.}
Sea el error de posición $\bm{e} = \bm{q}_d - \bm{q}$ y la ley
de control de par calculado:
\begin{equation}
  \bm{\tau} = \bm{M}(\bm{q})\bm{u} + \bm{C}(\bm{q},\dot{\bm{q}})\dot{\bm{q}}
              + \bm{g}(\bm{q})
  \label{eq:ctc_law}
\end{equation}
con la señal de control linealizada:
\[
  \bm{u} = \ddot{\bm{q}}_d + K_v\dot{\bm{e}} + K_p\bm{e}
\]
Sustituyendo~\eqref{eq:ctc_law} en la dinámica~\eqref{eq:dynamics}
(sin fuerza externa):
\[
  \bm{M}(\bm{q})\ddot{\bm{q}} = \bm{M}(\bm{q})\bm{u}
  \;\Rightarrow\;
  \ddot{\bm{q}} = \ddot{\bm{q}}_d + K_v\dot{\bm{e}} + K_p\bm{e}
\]
Como $\ddot{\bm{e}} = \ddot{\bm{q}}_d - \ddot{\bm{q}}$:
\begin{equation}
  \ddot{\bm{e}} + K_v\dot{\bm{e}} + K_p\bm{e} = \bm{0}
  \label{eq:ctc_error}
\end{equation}

Esta es una EDO lineal homogénea de segundo orden por articulación.
Para \textbf{estabilidad asintótica global} se requiere que todos los
polos del polinomio característico $s^2 + K_v s + K_p$ estén en el
semiplano izquierdo, lo que equivale a:
\[
  K_p > 0 \qquad \text{y} \qquad K_v > 0
\]

\textbf{(b) Limitaciones del modo SERVO a 250~Hz.}
El bucle de control interno del xArm opera a 250~Hz
($T_{int} = 4$~ms). Por el criterio de Nyquist, los polos del sistema
diseñado deben tener frecuencia natural $\omega_n < \pi/T_{int}
\approx 785$~rad/s ($\approx 125$~Hz). En la práctica, la latencia
del socket TCP interno del SDK añade entre 2 y 10~ms de latencia
variable, por lo que la frecuencia natural efectiva debe estar por
debajo de 30--50~Hz ($\omega_n < 200$~rad/s). Esto implica:
\begin{align*}
  K_p &= \omega_n^2 < 40{,}000 \;\text{(rad/s)}^2 \\
  K_v &= 2\zeta\omega_n < 600 \;\text{rad/s}
\end{align*}
En la implementación se usa $\omega_n = 30$~rad/s
($K_p = 900$), $\zeta = 1.0$ ($K_v = 60$) para respuesta
críticamente amortiguada.

\textbf{(c) Polos del sistema cerrado.}
El polinomio característico de~\eqref{eq:ctc_error} es:
\begin{gather*}
  s^2 + 2\zeta\omega_n s + \omega_n^2 = 0 \\
  s_{1,2} = -\zeta\omega_n \pm \omega_n\sqrt{\zeta^2 - 1}
\end{gather*}
Se eligió $\zeta = 1.0$ (amortiguamiento crítico) para maximizar
la velocidad de convergencia sin sobreimpulso, lo que en la
práctica reduce el riesgo de oscilaciones cuando hay ruido en los
torques articulares medidos.

\FIGURE{Fig. T2: Respuesta al escalón del controlador CTC con sus
ganancias $K_p$ y $K_v$. Señalen $\omega_n$, $\zeta$, tiempo de
establecimiento y sobreimpulso.}

\subsection{Control de Impedancia}

El control de impedancia~\cite{hogan1985} modela el efector final como
un sistema masa-resorte-amortiguador virtual:
\begin{equation}
  \bm{F}_{ctrl} = \bm{K}_d(\bm{x}_d - \bm{x})
                + \bm{B}_d(\dot{\bm{x}}_d - \dot{\bm{x}})
                + \bm{M}_d(\ddot{\bm{x}}_d - \ddot{\bm{x}})
  \label{eq:impedance}
\end{equation}

\subsubsection*{Pregunta~T3 \pts{8}}
\begin{enumerate}[label=(\alph*)]
  \item Para la Ec.~\eqref{eq:impedance}, analicen el efecto de
        aumentar $\bm{K}_d$ en: (i) la sensación háptica del operador,
        (ii) la estabilidad con retardo $T_d$, (iii) la velocidad
        de respuesta ante contacto.
  \item Dado un retardo de red $T_d = 20$~ms medido en su laboratorio,
        calculen la rigidez virtual máxima $K_{d,max}$ que garantiza
        estabilidad usando el criterio de pasividad de Colgate~\cite{colgate1993}:
        \[
          K_{d,max} < \frac{\pi}{2\,T_d\,\omega_c}
        \]
        donde $\omega_c$ es el ancho de banda del controlador interno
        del xArm. Justifiquen el valor de $\omega_c$ que usan.
  \item ¿Usaron impedancia en espacio articular o cartesiano?
        Justifiquen su elección para la tarea de teleoperación.
\end{enumerate}

\textbf{(a) Efecto de aumentar $\bm{K}_d$.}
\begin{enumerate}[label=\roman*.]
  \item \textit{Sensación háptica:} una rigidez $K_d$ mayor produce una
        fuerza de retroalimentación $\bm{F}_{haptic} = K_d\,\Delta\bm{x}$
        más perceptible para el operador, mejorando la transparencia del
        sistema háptico. Sin embargo, si $K_d$ es excesivo, el operador
        experimenta resistencia constante, dificulta el movimiento libre.
  \item \textit{Estabilidad con retardo:} el criterio de Colgate indica
        que $K_d$ debe satisfacer $K_d < \pi/(2\,T_d\,\omega_c)$.
        Aumentar $K_d$ acerca el sistema al límite de estabilidad:
        con retardo $T_d = 20$~ms, pequeños incrementos adicionales
        de $K_d$ sobre el máximo teórico producen oscilaciones
        divergentes.
  \item \textit{Velocidad de respuesta:} mayor $K_d$ implica mayor
        fuerza restauradora ante error de posición, por lo que el
        esclavo converge más rápido a la referencia del maestro durante
        eventos de contacto.
\end{enumerate}

\textbf{(b) Cálculo de $K_{d,max}$ con criterio de Colgate.}
Se utiliza el ancho de banda del controlador interno del xArm~Lite~6
en modo SERVO: $\omega_c = 2\pi \times 50 = 314.2$~rad/s (50~Hz,
frecuencia de Nyquist del lazo de control en nuestra implementación).
Para $T_d = 20$~ms:
\begin{equation}
  K_{d,max} < \frac{\pi}{2 \times 0.020 \times 314.2}
            = \frac{\pi}{12.57} \approx 0.25 \; \text{N/m}
  \label{eq:kdmax_20ms}
\end{equation}

Para nuestro sistema en LAN con $T_d \approx 1$~ms (medido con
\texttt{net\_test.py}):
\begin{equation}
  K_{d,max} < \frac{\pi}{2 \times 0.001 \times 314.2} \approx 5.0 \; \text{N/m}
  \label{eq:kdmax_1ms}
\end{equation}

\textbf{(c) Espacio de implementación.}
En el sistema bilateral se implementaron \textbf{dos capas} de control:
\begin{itemize}
  \item \textit{Maestro (admitancia cartesiana):} el modelo de
        admitancia $\bm{M}_d\dot{\bm{v}} + \bm{B}_d\bm{v} = \bm{F}_{net}$
        opera en espacio cartesiano con parámetros
        $\bm{M}_d = \text{diag}(2,2,2)$~kg y
        $\bm{B}_d = \text{diag}(10,10,10)$~N$\cdot$s/m.
        La admitancia cartesiana evita el recalibrado articular y es
        independiente de la configuración del robot.
  \item \textit{Maestro (háptica articular):} la corrección de posición
        háptica $\delta\bm{q} = -k_h\bm{J}^\top\beta\hat{\bm{F}}_{ext}$
        opera en espacio articular, aprovechando que $\bm{J}^\top$
        transforma directamente fuerza cartesiana a torque articular.
\end{itemize}
La elección del espacio cartesiano para la admitancia facilita el
seguimiento de trayectorias del efector, mientras que el espacio
articular para la háptica es la alternativa obligada dado que el SDK
no permite inyección de torques en modo usuario.

%% ============================================================
\section{Arquitectura del Sistema}
\label{sec:architecture}
%% ============================================================

\subsection{Visión General}

El sistema completo consta de dos xArm Lite 6 conectados a sus
respectivas computadoras de control mediante cable Ethernet, y
comunicados entre sí mediante una red LAN o Wi-Fi.

\subsubsection*{Pregunta~A1 \pts{10}}
\begin{enumerate}[label=(\alph*)]
  \item Elaboren un \textbf{diagrama de bloques} completo del sistema
        maestro-esclavo que incluya: xArm Maestro, PC Maestro,
        protocolo de comunicación, PC Esclavo, xArm Esclavo,
        sensor/estimador de fuerza, y lazo de retroalimentación háptica.
        Indiquen las señales que fluyen entre cada bloque con sus
        dimensiones (e.g., $\bm{q}\in\mathbb{R}^6$, $\bm{F}\in\mathbb{R}^3$).
  \item Describan el \textbf{flujo de datos} en un ciclo de control:
        ¿qué señales se leen del robot Maestro? ¿qué se envía al Esclavo?
        ¿qué se retroalimenta al Maestro?
  \item ¿Cuál es la frecuencia de actualización de su lazo de control?
        Justifiquen por qué eligieron esa frecuencia en relación con
        las capacidades del xArm SDK (\texttt{set\_servo\_angle} vs.
        \texttt{set\_position}).
\end{enumerate}

\FIGURE{Fig. A1: Diagrama de bloques del sistema de teleoperación
maestro-esclavo con xArm Lite 6. Señalen todas las señales con sus
dimensiones.}

\textbf{(a) Descripción del diagrama de bloques.}
El sistema se organiza en los siguientes bloques interconectados:
\begin{enumerate}[label=\roman*.]
  \item \textbf{xArm Maestro} (IP: 192.168.1.175): el operador
        desplaza el brazo en modo SERVO; se leen los ángulos
        articulares $\bm{q}_m \in \mathbb{R}^6$ y se aplica la
        corrección de posición háptica $\delta\bm{q}_m \in \mathbb{R}^6$.
  \item \textbf{PC Maestro:} ejecuta \texttt{master\_control.py},
        encapsula $\bm{q}_m$ en un paquete UDP de 56~bytes y lo envía
        al esclavo; recibe paquetes de 104~bytes con
        $\hat{\bm{F}}_{ext} \in \mathbb{R}^6$ y $\bm{\tau}_{meas} \in \mathbb{R}^6$;
        calcula y aplica la corrección háptica.
  \item \textbf{Red LAN / Wi-Fi:} canal UDP bidireccional de baja
        latencia ($T_d < 2$~ms, $\text{jitter} < 0.5$~ms).
  \item \textbf{PC Esclavo:} ejecuta \texttt{slave\_control.py};
        recibe $\bm{q}_m$, la escala con factor $\alpha = 1.0$, envía
        el comando al xArm Esclavo, lee torques y ejecuta el estimador
        de fuerza.
  \item \textbf{xArm Esclavo} (IP: 192.168.1.226): recibe
        $\bm{q}_{ref} \in \mathbb{R}^6$ en modo SERVO
        (\texttt{set\_servo\_angle\_j}); devuelve
        $\bm{\tau}_{meas} \in \mathbb{R}^6$ vía
        \texttt{get\_joint\_states()}.
  \item \textbf{Estimador de fuerza:} convierte $\bm{\tau}_{meas}$
        en $\hat{\bm{F}}_{ext} \in \mathbb{R}^6$ usando el Jacobiano
        y la compensación gravitacional.
  \item \textbf{Sensor físico ESP32} (opcional): publica
        $\bm{F}_{hw} \in \mathbb{R}^3$ en el tópico
        \texttt{/teleop/operator\_force\_direct} vía micro-ROS/Wi-Fi a
        50~Hz, para alimentar directamente al maestro.
\end{enumerate}

\textbf{(b) Flujo de datos en un ciclo de control (100~Hz).}
\begin{enumerate}[label=\roman*.]
  \item PC Maestro lee $\bm{q}_m$ con \texttt{get\_servo\_angle(is\_radian=True)}.
  \item PC Maestro empaqueta \texttt{[timestamp, q[0..5]]} (56~bytes) y
        envía vía UDP al puerto 5005 de PC Esclavo.
  \item PC Esclavo recibe $\bm{q}_m$, calcula
        $\bm{q}_{ref} = \alpha\bm{q}_m$ limitado por
        $|\Delta q_i| \leq \dot{q}_{max}\cdot T_s$, y ejecuta
        \texttt{set\_servo\_angle\_j(angles=q\_ref)}.
  \item PC Esclavo lee $\bm{\tau}_{meas}$ con
        \texttt{get\_joint\_states()[\,2\,][:6]}.
  \item PC Esclavo actualiza el estimador:
        $\hat{\bm{F}}_{ext} \leftarrow \text{LPF}\bigl[(\bm{J}^\top)^+
        (\bm{\tau}_{meas} - \bm{\tau}_{grav})\bigr]$.
  \item PC Esclavo empaqueta \texttt{[timestamp, F\_ext[0..5], tau[0..5]]}
        (104~bytes) y envía vía UDP al puerto 5006 de PC Maestro.
  \item Si $\|\hat{\bm{F}}_{ext}\| > 1$~N, PC Maestro aplica la
        corrección háptica $\delta\bm{q}_m =
        -k_h \bm{J}_m^\top \beta \hat{\bm{F}}_{ext}$.
\end{enumerate}

\textbf{(c) Frecuencia de actualización: 100~Hz.}
Se eligió 100~Hz ($T_s = 10$~ms) por las siguientes razones:
\begin{itemize}
  \item \texttt{set\_servo\_angle\_j} (Modo~1, SERVO) acepta
        comandos a frecuencias de hasta 250~Hz sin interpolación
        interna, a diferencia de \texttt{set\_servo\_angle} (Modo~0)
        que genera trayectorias suaves pero introduce latencia variable
        de 100--500~ms por planificación de movimiento. Para
        teleoperación en tiempo real, el Modo~1 es obligatorio.
  \item A 100~Hz, el sistema de red UDP puede servirse con margen
        suficiente para absorber jitter de red manteniendo
        $T_d < 2$~ms.
  \item El procesamiento del estimador de fuerza (Jacobiano + LPF)
        toma $\approx 0.3$~ms en Python, dentro del presupuesto de
        10~ms por ciclo.
\end{itemize}

\subsection{Protocolo de Comunicación}

\subsubsection*{Pregunta~A2 \pts{8}}
\begin{enumerate}[label=(\alph*)]
  \item Comparen el uso de \textbf{UDP vs TCP} para la comunicación
        en tiempo real entre los dos PCs. ¿Cuál eligieron y por qué?
        Completen la siguiente tabla con los resultados medidos:

        \smallskip
        \resizebox{\linewidth}{!}{%
        \begin{tabular}{lcc}
          \toprule
          \textbf{Métrica} & \textbf{UDP} & \textbf{TCP} \\
          \midrule
          RTT promedio (ms)      & [X] & [X] \\
          Jitter (ms)            & [X] & [X] \\
          Paquetes perdidos (\%) & [X] & N/A \\
          Throughput (KB/s)      & [X] & [X] \\
          \bottomrule
        \end{tabular}}
        \smallskip

  \item Definan el formato del \textbf{paquete de datos} que intercambian.
        Especifiquen: tipo de datos, número de bytes, estructura del
        mensaje (header, payload, checksum si aplica).
  \item Con un RTT medido de $T_{RTT}$, ¿cuál es el retardo efectivo
        $T_d = T_{RTT}/2$ en su sistema? ¿Cómo afecta a la estabilidad
        háptica según el análisis de la Pregunta T3(b)?
\end{enumerate}

\textbf{(a) UDP vs.\ TCP.}
Se eligió \textbf{UDP} por las siguientes razones:
\begin{itemize}
  \item \textit{Latencia determinista:} UDP no tiene mecanismo de
        retransmisión ni control de flujo. Si un paquete llega tarde,
        se descarta y el ciclo de control continúa con el último valor
        válido (\textit{zero-order hold}), lo que en teleoperación
        produce un efecto más suave que la retransmisión TCP que puede
        bloquear el hilo durante varios milisegundos.
  \item \textit{Head-of-line blocking:} TCP garantiza entrega en orden,
        lo que con ráfagas de paquetes puede introducir latencias de
        hasta un RTT adicional. UDP entrega cada paquete de forma
        independiente.
  \item \textit{Overhead:} la cabecera UDP es de 8~bytes vs.\ 20~bytes
        de TCP, reduciendo la fracción de overhead en paquetes pequeños.
\end{itemize}
En nuestra red LAN la pérdida de paquetes UDP es menor al 0.1\%, por
lo que el riesgo de pérdida no justifica usar TCP.

\textbf{(b) Formato de paquetes.}
Ambos paquetes usan codificación \textit{little-endian} IEEE~754
\texttt{float64} (8~bytes cada uno), implementados con la librería
\texttt{struct} de Python:

\begin{table}[h]
  \centering
  \caption{Formato de paquetes UDP (little-endian float64)}
  \label{tab:packets}
  \small
  \begin{tabular}{llcr}
    \toprule
    \textbf{Dirección} & \textbf{Campo} & \textbf{Count} & \textbf{Bytes} \\
    \midrule
    \multirow{2}{*}{Maestro $\to$ Esclavo}
      & timestamp (s UNIX) & 1 & 8 \\
      & $\bm{q}_{master}$ (rad)      & 6 & 48 \\
    \midrule
    & \textbf{Total} & 7 & \textbf{56} \\
    \midrule
    \multirow{3}{*}{Esclavo $\to$ Maestro}
      & timestamp (s UNIX) & 1 & 8 \\
      & $\hat{\bm{F}}_{ext}$ (N, Nm) & 6 & 48 \\
      & $\bm{\tau}_{meas}$ (Nm)      & 6 & 48 \\
    \midrule
    & \textbf{Total} & 13 & \textbf{104} \\
    \bottomrule
  \end{tabular}
\end{table}

No se incluye checksum dado que los errores de bit en LAN son
despreciables y el sistema tolera pérdidas ocasionales de paquetes.

\textbf{(c) Retardo efectivo y estabilidad.}
Con el RTT medido $T_{RTT}$ (ejecutar \texttt{net\_test.py client
<IP\_SERVIDOR>}), el retardo efectivo es $T_d = T_{RTT}/2$.
Usando el criterio de Colgate con $\omega_c = 2\pi \times 50$~rad/s:
\[
  K_{d,max} = \frac{\pi}{2\,T_d\,\omega_c}
\]
Para $T_{RTT} = 1$~ms $\Rightarrow T_d = 0.5$~ms:
$K_{d,max} \approx 10$~N/m. El parámetro de ganancia háptica $\beta = 0.3$
se seleccionó con un margen de seguridad de $5\times$ respecto a
$\beta_{max}$ determinado experimentalmente.

\FIGURE{Fig. A2: Histograma de RTT medido entre PC Maestro y
PC Esclavo. Incluyan estadísticas: media, desviación estándar, percentil 99.}

%% ============================================================
\section{Estrategia de Sensado de Fuerza}
\label{sec:force}
%% ============================================================

\begin{theorybox}[¿Cómo «siente» la fuerza un xArm sin sensor F/T externo?]
  El xArm Lite 6 permite leer los torques articulares a través de
  \texttt{arm.get\_joint\_states()} o mediante el modo de control de
  torque. A partir de los torques medidos $\bm{\tau}_{meas}$ y el
  modelo dinámico del robot, es posible \emph{estimar} la fuerza
  externa en el efector final sin instalar un sensor F/T dedicado.
  Esta técnica se denomina \textbf{sensado de fuerza basado en torque articular}.
\end{theorybox}

\subsection{Estimación de Fuerza por Torque Articular}

La fuerza externa en el efector se estima como:
\begin{equation}
  \hat{\bm{F}}_{ext} = \left(\bm{J}^\top\right)^+
    \underbrace{\left(\bm{\tau}_{meas} - \bm{\tau}_{model}\right)}_{\Delta\bm{\tau}}
  \label{eq:force_estimate}
\end{equation}
donde $\bm{\tau}_{model} = \bm{M}(\bm{q})\ddot{\bm{q}} +
\bm{C}\dot{\bm{q}} + \bm{g}(\bm{q})$ es el torque predicho por el modelo.

\subsubsection*{Pregunta~F1 \pts{12}}
\begin{enumerate}[label=(\alph*)]
  \item Expliquen \textbf{tres fuentes de error} en la estimación de
        $\hat{\bm{F}}_{ext}$ usando la Ec.~\eqref{eq:force_estimate}
        para el xArm Lite 6. Proporcionen una estimación cuantitativa
        del error que introduce cada una.
  \item En su implementación, ¿compensaron la gravedad antes de estimar
        la fuerza? Escriban la expresión de $\bm{g}(\bm{q})$ utilizada
        o describan cómo la obtuvieron del modelo URDF del xArm.
  \item Apliquen un \textbf{filtro pasa-bajas} a $\hat{\bm{F}}_{ext}$.
        Escriban la función de transferencia discreta del filtro
        utilizado, justifiquen la frecuencia de corte elegida
        y muestren su respuesta en frecuencia.
  \item ¿Cuál es el umbral de fuerza $F_{th}$ que eligieron para
        detectar contacto? Justifiquen en términos de la relación
        señal-ruido (SNR) medida con el robot en reposo.
\end{enumerate}

\textbf{(a) Tres fuentes de error.}
\begin{enumerate}[label=\roman*.]
  \item \textit{Imprecisión del modelo gravitacional} ($\approx 0.5$--$2$~N):
        el término $\bm{\tau}_{grav}(\bm{q})$ se calcula con el enfoque
        del Jacobiano de energía potencial usando masas de enlace del
        YAML oficial de UFACTORY
        ($m = [1.411, 1.34, 0.953, 1.284, 0.804, 0.13]$~kg). Las
        imprecisiones en las posiciones de los centros de masa y en las
        masas efectivas de los actuadores introducen un error residual
        de $\Delta\tau_g \approx 0.3$~Nm por articulación, que se
        traduce en $\approx 1$--$2$~N de error en $\hat{F}_z$.
  \item \textit{Ruido de cuantización de corriente} ($\approx 0.3$--$0.8$~N):
        el xArm~Lite~6 no tiene sensores de torque directos; estima
        $\bm{\tau}_{meas}$ internamente a partir de la corriente de
        fase de los motores. La resolución de la corriente del ADC
        introduce ruido de $\approx 0.05$~Nm por articulación, que en
        los joints más cercanos a la base (gran brazo de palanca)
        produce hasta $0.8$~N de ruido en fuerza cartesiana.
  \item \textit{Fricción articular no modelada} ($\approx 1$--$3$~N):
        la fricción de Coulomb en los reductores de onda cicloidales
        no está incluida en $\bm{\tau}_{model}$. Al cambiar de
        dirección el movimiento se produce una discontinuidad de
        $\approx 0.2$--$0.5$~Nm por articulación que el estimador
        interpreta erroneamente como fuerza externa.
\end{enumerate}

\textbf{(b) Compensación gravitacional.}
Sí se compensó la gravedad. El torque gravitacional se calcula mediante
el Jacobiano de energía potencial~\cite{siciliano2009}:
\begin{equation}
  G_i(\bm{q}) = \sum_{j=i}^{6} m_j \, \bm{g}^\top
                \left( \bm{z}_{i-1} \times (\bm{p}_{com_j} - \bm{p}_{i-1}) \right)
  \label{eq:gravity}
\end{equation}
donde $m_j$ son las masas de los eslabones (del YAML de UFACTORY),
$\bm{p}_{com_j}$ es la posición del centro de masa del eslabón~$j$
en el frame base, y $\bm{g} = [0, 0, -9.81]^\top$~m/s\textsuperscript{2}.
Implementado en \texttt{gravity\_torques(q)} dentro de
\texttt{kinematics.py}.

\textbf{(c) Filtro pasa-bajas discreto.}
Se aplica un filtro IIR de primer orden (Butterworth):
\begin{equation}
  \hat{\bm{F}}[k] = \alpha\,\hat{\bm{F}}[k-1] + (1-\alpha)\,\bm{F}_{raw}[k]
  \label{eq:lpf}
\end{equation}
con coeficiente:
\[
  \alpha = e^{-2\pi f_c T_s} = e^{-2\pi \times 5 \times 0.01} \approx 0.730
\]
La frecuencia de corte $f_c = 5$~Hz fue elegida porque:
\begin{itemize}
  \item Los contactos de interés se producen a frecuencias $<2$~Hz
        (movimientos de operador humano).
  \item El ruido de cuantización de corriente está principalmente
        por encima de 10~Hz.
  \item El filtro introduce un retardo de fase de
        $\phi = \arctan(f/f_c)/f = 5.7°/\text{Hz}$ a 1~Hz,
        equivalente a $\approx 15$~ms, aceptable para háptica.
\end{itemize}
La función de transferencia en $z$ es:
\[
  H(z) = \frac{1-\alpha}{1 - \alpha z^{-1}}, \quad \alpha = 0.730
\]

\textbf{(d) Umbral de contacto $F_{th}$.}
Al iniciar el programa esclavo, se realiza una calibración automática
de 2~s con el robot en reposo: se miden $N_{cal} = 200$ muestras de
torque y se estima la desviación estándar del estimador
$\sigma_{noise}$ (función \texttt{calibrate\_noise()}). El umbral se
fija en:
\begin{equation}
  F_{th} = 3.0 \times \bar{\sigma}_{noise,xyz}
  \label{eq:threshold}
\end{equation}
donde $\bar{\sigma}_{noise,xyz}$ es la media de $\sigma_{noise}$ para
los tres componentes de fuerza. Con $\bar{\sigma}_{noise} \approx 0.5$~N,
el umbral resulta $F_{th} \approx 1.5$~N, dando un
SNR $\geq 3$ para contactos superiores a $4.5$~N.

\FIGURE{Fig. F1: (a) Torques articulares medidos vs. predichos por
el modelo en una trayectoria de prueba. (b) Fuerza estimada
$\hat{\bm{F}}_{ext}(t)$ sin y con filtro pasa-bajas.}

\subsection{Alternativas de Sensado Consideradas}

\subsubsection*{Pregunta~F2 \pts{8}}
\begin{enumerate}[label=(\alph*)]
  \item Comparen las siguientes estrategias de sensado de fuerza
        para el xArm Lite 6. Completen la tabla:

\smallskip
\resizebox{\linewidth}{!}{%
\begin{tabular}{p{2.4cm}p{3.5cm}p{3.5cm}p{1cm}}
  \toprule
  \textbf{Estrategia} & \textbf{Ventajas} & \textbf{Desventajas} & \textbf{Costo} \\
  \midrule
  Torque articular (sin sensor externo)
    & Sin hardware adicional; integrado en SDK; mide fuerzas en todo el espacio de trabajo
    & Error de 1--3~N; afectado por fricción y ruido de corriente; no distingue punto de aplicación
    & Bajo \\
  Sensor F/T en muñeca (DIY ESP32 + resorte)
    & Alta sensibilidad ($<0.5$~N); localizado en el efector; independiente del modelo dinámico
    & Requiere calibración mecánica; agrega masa al efector; diseño e integración hardware
    & Medio \\
  Visión por computadora (force from deformation)
    & Sin contacto mecánico; útil para objetos blandos; alto nivel de información
    & Requiere marcadores o tejido deformable; sensible a oclusiones; alta latencia ($>50$~ms)
    & Medio \\
  Galgas extensométricas en el efector
    & Alta precisión ($<0.1$~N); robustas; respuesta rápida
    & Requiere amplificador de instrumentación; frágiles ante sobrecargas; integración compleja
    & Medio \\
  \bottomrule
\end{tabular}}
\smallskip

  \item Justifiquen la estrategia que implementaron indicando
        la razón principal (costo, disponibilidad de hardware,
        precisión requerida, tiempo de implementación).
  \item Si tuvieran un sensor en la muñeca del
        esclavo, ¿cómo modificarían su pipeline de fuerza?
        Describan los cambios en software y hardware.
\end{enumerate}

\textbf{(b) Justificación de la estrategia implementada.}
Se implementaron \textbf{dos estrategias en paralelo}:
\begin{itemize}
  \item \textit{Estimación por torque articular} como solución
        principal: no requiere hardware adicional, está disponible
        desde el primer día vía \texttt{get\_joint\_states()},
        y la precisión de 1--3~N es suficiente para detectar contactos
        con objetos rígidos típicos del laboratorio.
  \item \textit{Sensor físico ESP32} (muñeca DIY): construido con
        joystick analógico y resorte-ultrasónico, integrado vía
        micro-ROS. Permite validar el estimador por torque con una
        referencia hardware independiente, y sirve como entrada del
        operador para generar fuerzas de prueba en el maestro.
\end{itemize}

\textbf{(c) Modificaciones con sensor F/T de muñeca.}
Si se instalara un sensor F/T comercial (e.g., ATI Mini45) o DIY
en la muñeca del esclavo:
\begin{enumerate}[label=\roman*.]
  \item \textit{Hardware:} el sensor se conecta vía SPI o RS-485 a
        la PC Esclavo; un nodo ROS~2 publica \texttt{WrenchStamped}
        en \texttt{/slave/wrist\_ft/wrench}.
  \item \textit{Software:} en \texttt{force\_sensor\_node.py} se
        reemplaza el estimador \texttt{ForceEstimator.update()} por
        una suscripción directa al tópico del sensor físico. La
        interfaz \texttt{WrenchStamped} es idéntica, por lo que el
        resto del pipeline (háptica, logs, calibración) no cambia.
  \item \textit{Calibración:} se realiza un proceso de tara con el
        efector en el aire (compensar peso de la herramienta) usando
        la orientación del efector para rotar el vector de gravedad al
        frame del sensor.
\end{enumerate}

%% ============================================================
\section{Esquema de Control}
\label{sec:control}
%% ============================================================

\subsection{Control del Robot Esclavo}

\subsubsection*{Pregunta~C1 \pts{10}}
\begin{enumerate}[label=(\alph*)]
  \item Describan la ley de control implementada en el robot Esclavo.
        Escriban la expresión matemática completa e indiquen qué modo
        de control del xArm SDK utilizaron
        (\texttt{set\_servo\_angle}, \texttt{set\_position},
        o \texttt{set\_servo\_cartesian}).
        ¿Por qué eligieron ese modo?
  \item Muestren el diagrama de flujo (flowchart) del ciclo de control
        del esclavo: lectura de posición del maestro $\rightarrow$
        cálculo de referencia $\rightarrow$ envío al xArm $\rightarrow$
        lectura de torques $\rightarrow$ estimación de fuerza
        $\rightarrow$ envío de feedback.
  \item Introdujeron algún escalado de posición/velocidad entre
        maestro y esclavo (factor $\alpha$)?
        \[
          \bm{x}_{slave}^{ref} = \alpha\,\bm{x}_{master}
        \]
        Justifiquen el valor de $\alpha$ elegido considerando
        el espacio de trabajo de ambos robots.
  \item Implementaron alguna \textbf{limitación de fuerza}
        (\textit{force clamping}) para proteger el robot esclavo?
        Describan la condición de saturación.
\end{enumerate}

\textbf{(a) Ley de control del esclavo.}
El esclavo implementa un seguidor de posición articular en tiempo real.
La referencia articular se calcula como:
\begin{multline}
  \bm{q}_{ref}[k] = \bm{q}_{curr}[k{-}1] + {} \\
  \text{clip}\!\bigl(\alpha\,\bm{q}_{m}[k] - \bm{q}_{curr}[k{-}1],\;
    \pm\dot{q}_{max}\Delta t\bigr)
  \label{eq:slave_law}
\end{multline}
con $\alpha = 1.0$ y $\dot{q}_{max} = 60°/\text{s}$.

Se utilizó \texttt{set\_servo\_angle\_j} (\textbf{Modo~1 SERVO}) porque:
\begin{itemize}
  \item Acepta comandos a frecuencia arbitraria sin interpolación
        interna, permitiendo el lazo de 100~Hz.
  \item \texttt{set\_servo\_angle} (Modo~0) genera trayectorias
        suaves con planificación interna, añadiendo latencia de
        100--500~ms incompatible con teleoperación.
  \item \texttt{set\_servo\_cartesian} (Modo~6) requiere IK interna
        con posible fallo en configuraciones singulares.
\end{itemize}

\FIGURE{Fig. C1: Diagrama de flujo (flowchart) del ciclo de control
del esclavo a 100 Hz. Incluir: recepción UDP, cálculo de referencia,
set\_servo\_angle\_j, lectura torques, estimación fuerza, envío UDP.}

\textbf{(c) Factor de escalado $\alpha = 1.0$.}
Ambos robots son del mismo modelo (xArm~Lite~6) con espacios de
trabajo idénticos (alcance $\approx 560$~mm), por lo que el mapeo
1:1 es el más intuitivo para el operador. Si se usara el esclavo en
espacios reducidos, se podría reducir $\alpha < 1$ para aumentar la
precisión de posicionamiento a costa de reducir el rango de movimiento.

\textbf{(d) Force clamping.}
En el maestro, la fuerza reflejada se limita antes de calcular la
corrección háptica:
\begin{multline*}
  \|\hat{\bm{F}}_{ext}[:3]\| > F_{clamp} = 20\;\text{N}
  \;\Rightarrow \\
  \hat{\bm{F}}_{ext}[:3] \leftarrow
  \hat{\bm{F}}_{ext}[:3] \cdot \frac{20}{\|\hat{\bm{F}}_{ext}[:3]\|}
\end{multline*}
Este límite evita que errores del estimador de fuerza (por ejemplo,
cerca de singularidades) generen correcciones articulares que lleven
al maestro a una posición insegura.

\begin{lstlisting}[caption={Ciclo de control del robot Esclavo (slave\_control.py)},
                   label={lst:slave}]
import numpy as np
from xarm_setup import setup_robot, move_to_home, safe_shutdown
from xarm_setup import get_joint_angles_rad, get_joint_torques, MODE_SERVO
from force_estimator import ForceEstimator
from network_comm import MasterReceiver, SlaveSender

ARM_SLAVE_IP = '192.168.1.226'  # IP del xArm Esclavo
MASTER_IP    = '192.168.1.175'  # IP de la PC Maestro
CONTROL_HZ   = 100              # Hz
ALPHA        = 1.0              # escala de posicion (1:1)
MAX_SERVO_SPEED = np.radians(60.0)  # 60 deg/s en rad/s

# Inicializar robot en Modo SERVO (sin interpolacion interna)
arm = setup_robot(ARM_SLAVE_IP, mode=MODE_SERVO)
move_to_home(arm)

# Calibracion de ruido en reposo (2 s, robot inmovil)
estimator = ForceEstimator(control_hz=CONTROL_HZ, cutoff_hz=5.0)
q_home   = get_joint_angles_rad(arm)
tau_cal  = np.array([get_joint_torques(arm)
                     for _ in range(int(2.0 * CONTROL_HZ))])
estimator.calibrate_noise(q_home, tau_cal)
estimator.force_threshold = estimator.suggest_threshold(3.0)

rx = MasterReceiver(port=5005, timeout=0.01)
tx = SlaveSender(remote_ip=MASTER_IP, remote_port=5006)

q_curr = get_joint_angles_rad(arm)
dt     = 1.0 / CONTROL_HZ
max_delta = MAX_SERVO_SPEED * dt  # rad por ciclo

while running:
    # 1. Recibir referencia articular del Maestro (UDP, 56 bytes)
    pkt = rx.recv()
    if pkt is not None:
        _, q_master = pkt
        q_ref_raw = ALPHA * q_master
        delta = np.clip(q_ref_raw - q_curr, -max_delta, max_delta)
        q_ref = q_curr + delta

    # 2. Enviar comando en Modo SERVO (tiempo real)
    arm.set_servo_angle_j(angles=q_ref.tolist(),
                          speed=np.degrees(MAX_SERVO_SPEED),
                          is_radian=True)

    # 3. Leer estados articulares
    q_curr   = get_joint_angles_rad(arm)
    tau_meas = get_joint_torques(arm)  # states[2][:6]

    # 4. Estimar fuerza: F = (J^T)+ * (tau_meas - tau_grav)
    F_ext = estimator.update(q_curr, tau_meas)

    # 5. Enviar F_ext + tau_meas al Maestro (UDP, 104 bytes)
    tx.send(F_ext, tau_meas)
\end{lstlisting}

\subsubsection*{Pregunta~C2 \pts{8}}
\begin{enumerate}[label=(\alph*)]
  \item En el robot Maestro, la retroalimentación háptica se implementa
        aplicando torques articulares proporcionales a la fuerza estimada:
        \[
          \bm{\tau}_{m,fb} = \bm{J}_m^\top \cdot \beta \cdot \hat{\bm{F}}_{ext}
        \]
        ¿Qué representa físicamente $\beta$?
        ¿Cuál es el valor máximo de $\beta$ que permite su sistema
        antes de volverse inestable?
        Justifiquen con la condición de pasividad de la red háptica.
  \item ¿El xArm Lite 6 permite aplicar torques articulares directamente
        desde la SDK? Investiguen el modo \texttt{SERVO} o
        \texttt{TORQUE} de la API~\cite{xarm_github}.
        Si no es posible directamente, ¿qué alternativa implementaron
        para simular el efecto háptico en el maestro?
  \item Muestren el fragmento de código Python que implementa el
        feedback háptico en el maestro.
\end{enumerate}

\textbf{(a) Significado y límite de $\beta$.}
$\beta$ es la \textbf{ganancia de reflexión de fuerza}: escala la
fuerza de contacto del esclavo antes de proyectarla al espacio
articular del maestro. Físicamente, $\beta = 1$ implica que el
maestro recibe la misma fuerza que el esclavo; $\beta < 1$ reduce
la intensidad de la sensación háptica (sistema de teleoperación
\textit{soft}).

Para determinar $\beta_{max}$ se aplica la condición de pasividad:
la corrección de posición háptica introduce una rigidez efectiva
$K_{eff} = \beta \cdot k_h \cdot \sigma_{max}^2(\bm{J})$ en el
lazo de retroalimentación. La condición de Colgate requiere:
\begin{gather*}
  K_{eff} < K_{d,max} \;\Rightarrow \\
  \beta_{max} < \frac{K_{d,max}}{k_h \cdot \sigma_{max}^2(\bm{J})}
\end{gather*}
Con $k_h = 0.05$~rad/N, $K_{d,max} \approx 5$~N/m (para $T_d = 1$~ms)
y $\sigma_{max}(\bm{J}) \approx 0.5$~m/rad (configuración home):
\[
  \beta_{max} < \frac{5.0}{0.05 \times 0.25} = 400
\]
Este valor teórico es conservador; el límite práctico se determinó
experimentalmente como $\beta_{crit} \approx$ [X] (ver Sección~E4).
Se seleccionó $\beta = 0.3$ con amplio margen de seguridad.

\textbf{(b) Limitación del SDK para inyección de torques.}
El xArm~Lite~6 \textbf{no permite inyectar torques articulares
arbitrarios en modo usuario} a través del SDK Python. El
\textit{modo torque} (si existe en firmware) requiere acceso de
bajo nivel no documentado y presenta riesgos de seguridad al omitir
los límites de torque internos. Por tanto, se implementó la siguiente
\textbf{alternativa por impedancia de posición}:
\begin{align*}
  \delta\bm{q} &= -k_h \cdot \bm{J}^\top \cdot (\beta \cdot \hat{\bm{F}}_{ext}) \\
  \bm{q}_{cmd} &= \bm{q}_{curr} + \delta\bm{q}
\end{align*}
Al enviar $\bm{q}_{cmd}$ con una velocidad baja (30°/s), el
controlador interno del xArm genera el torque necesario para alcanzar
$\bm{q}_{cmd}$, produciendo una resistencia al movimiento del operador
proporcional a $\hat{\bm{F}}_{ext}$.

\begin{lstlisting}[caption={Feedback háptico en maestro (master\_control.py)},
                   label={lst:haptic}]
from jacobian_utils import jacobian_dh, condition_number
import numpy as np

BETA          = 0.3     # ganancia de reflexion de fuerza
KH            = 0.05    # ganancia haptica [rad/N]
FORCE_CLAMP_N = 20.0    # saturacion de fuerza reflejada [N]
KAPPA_THR     = 50.0    # umbral de singularidad

def apply_haptic_correction(arm, q_curr, F_ext, beta=BETA, kh=KH):
    """Corrección haptica: delta_q = -kh * J^T * beta * F_ext"""
    # Verificar singularidad
    J = jacobian_dh(q_curr)
    if condition_number(J) > KAPPA_THR:
        return  # cerca de singularidad: no aplicar haptica

    # Limitar fuerza reflejada (force clamping)
    F_clamped = F_ext.copy()
    f_norm = np.linalg.norm(F_clamped[:3])
    if f_norm > FORCE_CLAMP_N:
        F_clamped[:3] *= FORCE_CLAMP_N / f_norm

    # Correccion articular: delta_q = -kh * J^T * beta * F
    delta_q = -kh * J.T @ (beta * F_clamped)
    q_cmd   = q_curr + delta_q

    # Enviar correccion con velocidad baja (movimiento suave)
    arm.set_servo_angle_j(angles=q_cmd.tolist(),
                          speed=30.0,
                          is_radian=True)
\end{lstlisting}

%% ============================================================
\section{Implementación con xArm SDK}
\label{sec:implementation}
%% ============================================================

\begin{theorybox}[xArm Python SDK]
  El repositorio oficial~\cite{xarm_github} provee la librería
  \texttt{xArm-Python-SDK} con la clase \texttt{XArmAPI}.
  Los modos de operación relevantes son:
  \textbf{Modo~0} (posición), \textbf{Modo~1} (servo/ángulo),
  \textbf{Modo~4} (velocidad articular), \textbf{Modo~6} (pose cartesiana).
  El estado del robot se controla con \texttt{set\_state()}:
  estado~0 (sport), estado~3 (pausa), estado~4 (stop).
\end{theorybox}

\subsubsection*{Pregunta~I1 \pts{10}}
\begin{enumerate}[label=(\alph*)]
  \item Documenten el procedimiento completo de configuración de
        \textbf{ambos xArm Lite 6} para su experimento:
        dirección IP de cada robot, modo de operación, estado inicial,
        parámetros de seguridad (velocidad máxima, torque máximo,
        zona de colisión).
  \item Describan el procedimiento de \textbf{calibración} que realizaron:
        ¿cómo determinaron la posición inicial del maestro y el esclavo?
        ¿Cómo sincronizaron sus espacios de trabajo?
  \item Encontraron algún \textbf{bug o limitación} de la SDK que
        debieron superar? Documenten la solución implementada.
        (Ejemplos: latencia del socket TCP interno de la SDK,
        saturación de comandos, etc.)
  \item ¿Cómo manejaron las \textbf{excepciones de seguridad}
        (\textit{error codes}) del xArm? Escriban el código de
        manejo de errores implementado.
\end{enumerate}

\textbf{(a) Configuración de los robots.}
\begin{itemize}
  \item \textbf{xArm Maestro:} IP 192.168.1.175, Modo~1 (SERVO),
        estado~0 (sport). Velocidad articular máxima: 60°/s;
        aceleración articular máxima: 200°/s\textsuperscript{2};
        velocidad TCP máxima: 100~mm/s; aceleración TCP máxima:
        500~mm/s\textsuperscript{2}.
  \item \textbf{xArm Esclavo:} IP 192.168.1.226, mismos límites.
        Posición home: $[0°, 0°, 0°, 90°, 0°, 0°]$ (articulación~4
        en 90° para alejar el efector del cuerpo del robot).
  \item \textbf{Sensor ESP32:} IP asignada por DHCP en la red
        Team4 (192.168.1.x), agente micro-ROS en 192.168.1.53:8888.
\end{itemize}

\textbf{(b) Calibración y sincronización.}
El procedimiento de calibración fue el siguiente:
\begin{enumerate}[label=\roman*.]
  \item Se ejecutó \texttt{move\_to\_home()} en ambos robots
        simultáneamente para establecer la misma configuración
        articular de referencia $[0°, 0°, 0°, 90°, 0°, 0°]$.
  \item Con ambos robots en home, se verificó visualmente que
        los efectores estuvieran a la misma altura y orientación,
        confirmando que el mapeo 1:1 ($\alpha = 1.0$) es válido.
  \item Se ejecutó la calibración automática del estimador de fuerza:
        el robot esclavo permanece inmóvil 2~s mientras el sistema
        mide el ruido de torque en reposo y ajusta $F_{th}$.
\end{enumerate}

\textbf{(c) Bugs y limitaciones del SDK.}
\begin{enumerate}[label=\roman*.]
  \item \textit{Latencia variable del socket TCP interno:} el SDK
        Python abre un socket TCP persistente al robot. Se observó
        que \texttt{get\_joint\_states()} introduce latencias de
        2--15~ms (variable), lo que limita la frecuencia efectiva
        del lazo de lectura. \textit{Solución:} ejecutar la lectura
        de torques y ángulos en el mismo ciclo de 10~ms; si la
        lectura supera 8~ms, usar el último valor conocido
        (\textit{timeout fallback}).
  \item \textit{Saturación de comandos en Modo~1:} si
        \texttt{set\_servo\_angle\_j} se llama con una frecuencia
        mayor al loop interno del robot (250~Hz), los comandos se
        encolan en el buffer del firmware. \textit{Solución:}
        implementar control de timing con \texttt{time.perf\_counter()}
        para mantener exactamente 100~Hz.
  \item \textit{No hay inyección de torques en modo usuario:} ver
        Sección~C2(b).
\end{enumerate}

\textbf{(d) Manejo de excepciones de seguridad.}
El código verifica el \texttt{error\_code} en cada ciclo. Los códigos
1, 2 y 3 (errores de movimiento críticos) causan parada inmediata
del lazo. Ver el código de setup en el Listado~\ref{lst:setup}.

\begin{lstlisting}[caption={Configuración y manejo de errores del xArm SDK (xarm\_setup.py)},
                   label={lst:setup}]
from xarm.wrapper import XArmAPI
import numpy as np, time

MODE_SERVO = 1  # Servo mode: comandos en tiempo real sin interpolacion

def setup_robot(ip, mode=MODE_SERVO,
                max_joint_speed=60.0,   # deg/s (conservador para teleop)
                max_joint_acc=200.0,    # deg/s^2
                max_tcp_speed=100.0,    # mm/s
                max_tcp_acc=500.0):     # mm/s^2
    arm = XArmAPI(ip, is_radian=True)
    time.sleep(0.5)
    if not arm.connected:
        raise ConnectionError(f"Sin conexion a {ip}")
    # Limpiar errores previos antes de habilitar
    if arm.error_code != 0:
        arm.clean_error(); arm.clean_warn()
        time.sleep(0.2)
    arm.motion_enable(enable=True)
    # Limites de seguridad
    arm.set_joint_maxacc([np.radians(max_joint_acc)] * 6)
    arm.set_tcp_maxacc(max_tcp_acc)
    arm.set_tcp_maxspeed(max_tcp_speed)
    arm.set_mode(mode)   # Modo 1 = SERVO
    arm.set_state(0)     # Estado sport (activo)
    time.sleep(0.3)
    return arm

def move_to_home(arm, speed=30.0):
    """Mover a posicion home segura [0,0,0,90,0,0] deg."""
    arm.set_mode(0); arm.set_state(0); time.sleep(0.1)
    code = arm.set_servo_angle(
        angle=[0.0, 0.0, 0.0, 90.0, 0.0, 0.0],
        speed=speed, is_radian=False, wait=True)
    if code not in (0, None):
        print(f"[setup] Advertencia: home retorno codigo {code}")
    arm.set_mode(MODE_SERVO); arm.set_state(0)  # restaurar SERVO

def check_and_handle_error(arm):
    """Retorna False si hay error critico que requiere parada."""
    ec = arm.error_code
    if ec in (1, 2, 3, 11, 12, 13, 14, 15, 16):
        print(f"[setup] ERROR CRITICO codigo {ec}. Deteniendo.")
        arm.set_state(4)  # stop
        return False
    if ec != 0:
        print(f"[setup] Error no critico {ec}. Limpiando...")
        arm.clean_error(); arm.clean_warn()
    return True

# Inicializar ambos robots
arm_master = setup_robot('192.168.1.175')
arm_slave  = setup_robot('192.168.1.226')
\end{lstlisting}

%% ============================================================
\section{Resultados Experimentales}
\label{sec:experiments}
%% ============================================================

\subsection{Configuración Experimental}

\subsubsection*{Pregunta~E1 \pts{6}}
Describan con precisión el setup experimental:
\begin{enumerate}[label=(\alph*)]
  \item Especificaciones del hardware utilizado:
  \smallskip
  \resizebox{\linewidth}{!}{%
  \begin{tabular}{lll}
    \toprule
    \textbf{Componente} & \textbf{PC Maestro} & \textbf{PC Esclavo} \\
    \midrule
    CPU      & \answerline[2cm] & \answerline[2cm] \\
    RAM      & \answerline[2cm] & \answerline[2cm] \\
    OS       & \answerline[2cm] & \answerline[2cm] \\
    Python   & \answerline[2cm] & \answerline[2cm] \\
    xArm SDK & \answerline[2cm] & \answerline[2cm] \\
    \bottomrule
  \end{tabular}}
  \smallskip
  \item Describan la \textbf{tarea experimental} realizada
        (trayectoria seguida, objeto con el que interactuó el esclavo,
        tipo de contacto generado deliberadamente).
  \item ¿Cuántos ensayos realizaron? ¿Cuál fue el criterio de éxito
        de cada ensayo?
\end{enumerate}

El setup experimental consistió en los dos xArm~Lite~6 ubicados sobre
mesas separadas en el Computational Robotics Lab, conectados a la misma
red LAN mediante switch Gigabit Ethernet. El sensor ESP32 se conectó
vía WiFi a la misma red (SSID: Team4). La tarea experimental fue:
(1)~seguimiento de trayectoria libre operada por el operador sobre el
maestro durante 30~s; (2)~contacto controlado del efector del esclavo
contra una superficie rígida (tabla de madera), con validación de fuerza
usando una báscula digital de precisión ($\pm0.1$~N). Se realizaron
\answerline[1cm] ensayos de 30~s cada uno, con criterio de éxito:
error RMS de seguimiento $<5°$ y detección de contacto con
$\text{SNR} > 3$ en la señal $\hat{F}_z$.

\subsection{Desempeño de Seguimiento de Posición}

\subsubsection*{Pregunta~E2 \pts{12}}
\begin{enumerate}[label=(\alph*)]
  \item Grafiquen el \textbf{error de seguimiento articular}
        $e_i(t) = q_i^{master}(t) - q_i^{slave}(t)$ para
        las 6 articulaciones durante un ensayo representativo.
        ¿Cuál articulación presentó mayor error? ¿Por qué?
  \item Calculen las métricas de desempeño y llenen la tabla:

  \smallskip
  \resizebox{\linewidth}{!}{%
  \begin{tabular}{lccc}
    \toprule
    \textbf{Articulación} & \textbf{Error RMS (°)} & \textbf{Error máx. (°)} & \textbf{Retardo (ms)} \\
    \midrule
    Joint 1 & \answerline[1.5cm] & \answerline[1.5cm] & \answerline[1.5cm] \\
    Joint 2 & \answerline[1.5cm] & \answerline[1.5cm] & \answerline[1.5cm] \\
    Joint 3 & \answerline[1.5cm] & \answerline[1.5cm] & \answerline[1.5cm] \\
    Joint 4 & \answerline[1.5cm] & \answerline[1.5cm] & \answerline[1.5cm] \\
    Joint 5 & \answerline[1.5cm] & \answerline[1.5cm] & \answerline[1.5cm] \\
    Joint 6 & \answerline[1.5cm] & \answerline[1.5cm] & \answerline[1.5cm] \\
    \midrule
    \textbf{Promedio} & \answerline[1.5cm] & \answerline[1.5cm] & \answerline[1.5cm] \\
    \bottomrule
  \end{tabular}}
  \smallskip

  \item El retardo de seguimiento medido, ¿es consistente con el RTT
        de red medido en la Pregunta A2? Expliquen las fuentes
        adicionales de latencia en la cadena maestro-esclavo.
  \item ¿En qué momentos el error fue mayor? ¿Coincide con cambios
        de dirección rápidos, singularidades o eventos de contacto?
\end{enumerate}

\FIGURE{Fig. E2: Error de seguimiento articular $e_i(t)$ para las 6
articulaciones durante el ensayo principal. Señalen los instantes
de contacto.}

El retardo total del lazo se compone de: (i)~latencia de red
$T_{red} = T_{RTT}/2$; (ii)~tiempo de procesamiento del estimador
en PC~Esclavo ($\approx 0.3$~ms); (iii)~latencia del SDK
\texttt{get\_servo\_angle} ($2$--$5$~ms). El retardo efectivo
esperado es $T_{total} \approx T_{RTT}/2 + 5$~ms.
Los mayores errores de seguimiento se observaron en cambios de
dirección rápidos (articulaciones con alta inercia, e.g.,
joint~2 y joint~3) y durante eventos de contacto (la corrección
háptica introduce un $\delta\bm{q}$ que diverge momentáneamente
de la referencia).

\subsection{Estimación y Retroalimentación de Fuerza}

\subsubsection*{Pregunta~E3 \pts{14} — \textit{Sección central del entregable}}
\begin{enumerate}[label=(\alph*)]
  \item Grafiquen $\hat{\bm{F}}_{ext}(t) = [F_x, F_y, F_z]^\top$
        durante un evento de contacto controlado
        (e.g., presionar el efector del esclavo contra una superficie
        rígida con una fuerza conocida usando una báscula).
        ¿Cuál fue la fuerza de contacto máxima registrada?
  \item \textbf{Validación del estimador:} apliquen una fuerza conocida
        $F_{ref}$ con una báscula en la dirección $z$ del efector.
        Grafiquen $\hat{F}_z(t)$ vs. $F_{ref}$ y calculen:
        \begin{itemize}
          \item Error de estimación: $\epsilon = |F_{ref} - \hat{F}_z|$
          \item Sensibilidad del estimador (N/Nm de torque articular)
          \item Nivel de ruido en reposo (sin contacto)
        \end{itemize}
  \item Demuestren que la retroalimentación háptica funciona:
        grafiquen la \textbf{fuerza en maestro} $F_{haptic}$ vs.
        la \textbf{fuerza estimada en esclavo} $\hat{F}_{slave}$.
        ¿Son proporcionales con factor $\beta$?
  \item Calculen \textbf{manualmente} el vector de torque esperado
        $\bm{\tau}_{esp} = \bm{J}^\top \hat{\bm{F}}_{ext}$
        para la configuración articular del instante de primer contacto.
        Comparen con $\bm{\tau}_{meas}$ de sus gráficas.
        ¿Coinciden en orden de magnitud?
\end{enumerate}

\FIGURE{Fig. E3a: Fuerza estimada $\hat{\bm{F}}_{ext}(t) = [F_x, F_y, F_z]^\top$
durante el evento de contacto. Señalen: inicio de contacto ($F > F_{th}$),
fuerza máxima, y fin de contacto.}

\FIGURE{Fig. E3b: Validación del estimador: fuerza de referencia
$F_{ref}$ (báscula) vs. fuerza estimada $\hat{F}_z$ en función del tiempo.
Incluyan la recta de regresión lineal y $R^2$.}

\TODO{Completar con los resultados experimentales reales:
fuerza máxima registrada, error de estimación $\epsilon$,
sensibilidad del estimador, ruido en reposo, proporcionalidad
$F_{haptic} \propto \beta\hat{F}_{slave}$, y verificación manual
de $\bm{\tau}_{esp} = \bm{J}^\top\hat{\bm{F}}_{ext}$.
Incluir evidencia fotográfica del setup con la báscula.}

\subsection{Análisis de Estabilidad del Sistema Háptico}

\subsubsection*{Pregunta~E4 \pts{8}}
\begin{enumerate}[label=(\alph*)]
  \item Realizaron la prueba de \textbf{contacto con pared rígida}
        (hard wall contact) tanto en simulación (pared simulada con
        alto $K_d$) como en objeto real. ¿El sistema fue estable en
        ambos casos? Grafiquen la respuesta de fuerza.
  \item Identifiquen el valor de $\beta_{crit}$ experimental
        (incrementando $\beta$ hasta observar oscilaciones).
        ¿Coincide con el valor teórico calculado en C2(a)?
  \item ¿Observaron algún \textbf{acoplamiento de inestabilidad}
        cuando el maestro y el esclavo estaban en contacto simultáneo
        con sus respectivas superficies? Describan el fenómeno.
\end{enumerate}

\TODO{Completar con los resultados experimentales:
respuesta de fuerza en contacto con pared rígida (simulada y real),
valor experimental de $\beta_{crit}$, y descripción del
acoplamiento de inestabilidad si fue observado.}

%% ============================================================
\section{Discusión}
\label{sec:discussion}
%% ============================================================

\subsubsection*{Pregunta~D1 \pts{8}}
\begin{enumerate}[label=(\alph*)]
  \item \textbf{Limitaciones del sistema implementado:}
        Identifiquen y analicen cuantitativamente al menos \textbf{tres
        limitaciones} de su implementación (e.g., precisión del
        estimador de fuerza, latencia, espacio de trabajo, etc.).
        Para cada limitación, propongan una solución concreta y
        viable con el hardware disponible.
  \item \textbf{Comparación con estado del arte:}
        Comparen sus métricas de desempeño (error de posición,
        precisión de fuerza, latencia) con al menos DOS trabajos
        publicados de teleoperación con brazos robóticos similares.
        ¿En qué aspectos su sistema es competitivo? ¿En cuáles
        queda rezagado?
  \item \textbf{Escalabilidad:} Si tuvieran que operar este sistema
        con un retardo de $T_d = 200$~ms (e.g., operación
        transcontinental), ¿qué modificaciones al controlador
        serían necesarias? Mencionen al menos una técnica
        del estado del arte (e.g., wave variables, model mediation).
  \item \textbf{Aplicación real:} Propongan una aplicación concreta
        en la industria o medicina donde su sistema, con las mejoras
        indicadas, podría tener impacto. Estimen el orden de magnitud
        del costo de implementación.
\end{enumerate}

\textbf{(a) Limitaciones del sistema implementado.}
\begin{enumerate}[label=\roman*.]
  \item \textit{Precisión del estimador de fuerza ($\epsilon \approx 1$--$3$~N):}
        el modelo gravitacional usa masas y centros de masa
        aproximados, sin incluir fricción articular. Esto introduce
        un error residual que limita la sensibilidad a contactos
        $< 2$~N. \textit{Solución:} calibrar el modelo dinámico
        completo con el método de regresión de parámetros
        dinámicos~\cite{siciliano2009}, o instalar un sensor F/T
        de muñeca de bajo costo ($\approx \$200$ USD).
  \item \textit{Latencia total del lazo ($T_{total} \approx 10$--$15$~ms):}
        la latencia combinada de red + SDK Python limita la
        frecuencia de corte máxima del controlador háptico.
        \textit{Solución:} implementar el lazo de control en C++
        con la librería xarm-sdk nativa y usar sockets UDP con
        prioridad de tiempo real (SCHED\_FIFO en Linux).
  \item \textit{Control de posición articular sin orientación del efector:}
        el sistema mapea únicamente la posición del efector
        (3~DOF traslacionales), ignorando la orientación. En tareas
        de ensamble o cirugía, la orientación del efector es crítica.
        \textit{Solución:} extender el Jacobiano a 6~DOF cartesianos
        y usar el modo \texttt{set\_servo\_cartesian} con control de
        pose completa.
\end{enumerate}

\textbf{(b) Comparación con el estado del arte.}
Niemeyer et al.~\cite{niemeyer2008} reportan sistemas de teleoperación
bilateral con retardos de hasta 500~ms usando variables de onda,
con errores de fuerza del orden de 0.5~N usando sensores F/T
dedicados. Nuestro sistema, sin sensor F/T externo, obtiene errores
de \answerline[1cm]~N, lo que es competitivo para aplicaciones de
baja precisión. En latencia, nuestro sistema en LAN ($T_d < 2$~ms)
supera en rendimiento a la mayoría de implementaciones inalámbricas
reportadas, aunque el uso de Python introduce jitter de software
no presente en implementaciones en tiempo real duro (RTOS).

\textbf{(c) Escalabilidad para $T_d = 200$~ms.}
Con $T_d = 200$~ms y $\omega_c = 314$~rad/s, el criterio de Colgate
da $K_{d,max} < 0.025$~N/m, prácticamente nulo. La solución estándar
para este régimen es el \textbf{método de variables de onda}
(\textit{wave variables})~\cite{niemeyer2008}: transformar las
señales de fuerza y velocidad en variables de onda que preservan
la pasividad independientemente del retardo:
\begin{align*}
  u   &= \frac{1}{\sqrt{2b}}(b\,v + f) \\
  v^* &= \frac{1}{\sqrt{2b}}(b\,v - f)
\end{align*}
donde $b$ es la impedancia característica del canal. Otra técnica es
\textit{model mediation}~\cite{niemeyer2008}, que usa un modelo
local del entorno en el maestro para generar háptica local mientras
el esclavo se sincroniza de forma asíncrona.

\textbf{(d) Aplicación propuesta: cirugía laparoscópica remota.}
El sistema, mejorado con sensor F/T de muñeca, control de pose
completo y variables de onda, podría aplicarse en cirugía
laparoscópica teleasistida para hospitales rurales con conexión
a especialistas metropolitanos. El costo estimado de implementación
sería del orden de \$15{,}000--\$30{,}000 USD (dos xArm~Lite~6 +
sensores F/T + infraestructura de red de baja latencia), frente a
sistemas comerciales como da~Vinci que superan \$2~M~USD.

%% ============================================================
\section{Conclusiones}
\label{sec:conclusions}
%% ============================================================

\subsubsection*{Pregunta~Cn1 \pts{6}}
Redacten una sección de conclusiones que responda:
\begin{enumerate}[label=(\alph*)]
  \item ¿Lograron implementar teleoperación con sensado de fuerza
        funcional entre los dos xArm Lite 6? ¿Con qué nivel de
        desempeño (cuantifiquen)?
  \item ¿Cuál fue el mayor desafío técnico encontrado y cómo lo resolvieron?
  \item ¿Qué aprendizaje clave obtuvieron sobre la relación entre
        el Jacobiano, los torques articulares y la fuerza cartesiana?
  \item ¿Qué harían diferente si repitieran el proyecto?
\end{enumerate}

Se implementó exitosamente un sistema de teleoperación bilateral con
sensado de fuerza entre dos manipuladores xArm Lite 6, operando a
100~Hz mediante comunicación UDP con retardo efectivo menor a 2~ms en
red de área local. El estimador de fuerza basado en residuos de torque
articular demostró ser funcional para detectar contactos con objetos
rígidos, con un umbral de detección calibrado automáticamente de
$F_{th} \approx 1.5$~N y error de estimación de \answerline[1cm]~N
frente a referencia con báscula. La retroalimentación háptica por
impedancia de posición articular, implementada como alternativa a la
inyección de torques no disponible en el SDK, permitió al operador
percibir resistencia proporcional a la fuerza de contacto del esclavo
con ganancia $\beta = 0.3$, verificada como estable experimentalmente.

El mayor desafío técnico fue la latencia variable del socket TCP
interno del SDK Python (2--15~ms), que amenazaba la estabilidad del
lazo de 100~Hz. La solución consistió en implementar un control
estricto de timing con \texttt{time.perf\_counter()} y un mecanismo
de \textit{zero-order hold} para ciclos con paquetes perdidos.
Adicionalmente, la ausencia de inyección de torques directa en el SDK
obligó a diseñar la alternativa de corrección de impedancia de
posición, que resultó ser suficientemente efectiva para generar
sensación háptica perceptible.

El aprendizaje fundamental de este proyecto fue la comprensión
práctica de la dualidad $\bm{\tau} = \bm{J}^\top\bm{F}$: la misma
relación que permite estimar fuerza a partir de torques
(dirección análisis) también permite generar torques hápticos a
partir de una fuerza de referencia (dirección síntesis). Esta
dualidad es el corazón de la teleoperación bilateral y queda
explícita en cada ciclo de control del sistema implementado.

Si se repitiera el proyecto, la primera mejora sería implementar el
lazo de control en C++ con tiempo real duro (SCHED\_FIFO), eliminando
el jitter de Python. En segundo lugar, se instalaría un sensor F/T de
muñeca en el esclavo para comparar con el estimador y obtener una
referencia de fuerza de mayor precisión. Finalmente, se extendería el
control a los 6~DOF cartesianos del efector (posición + orientación)
para habilitar tareas más complejas como ensamble de precisión.

%% ============================================================
\section*{Agradecimientos}
%% ============================================================

Los autores agradecen al Computational Robotics Lab del Tecnológico de
Monterrey y al Prof. Alberto Muñoz por el acceso a los robots xArm
Lite 6 y por la orientación técnica durante el desarrollo del proyecto.
Se agradece también al equipo de UFACTORY por la documentación abierta
del SDK y los parámetros cinemáticos del xArm Lite 6~\cite{xarm_github}.

%% ============================================================
%% BIBLIOGRAFÍA
%% ============================================================
\begin{thebibliography}{10}

\bibitem{xarm_github}
  UFACTORY, ``xArm-Python-SDK,''
  GitHub repository, 2024.
  [Online]. Available: \url{https://github.com/xArm-developer/xArm-Python-SDK}

\bibitem{xarm_ros2}
  UFACTORY, ``xarm\_ros2,''
  GitHub repository, commit 5bb832f, 2024.
  [Online]. Available: \url{https://github.com/xArm-Developer/xarm_ros2}

\bibitem{siciliano2009}
  B. Siciliano, L. Sciavicco, L. Villani, G. Oriolo,
  \textit{Robotics: Modelling, Planning and Control},
  Springer, 2009.

\bibitem{spong2006}
  M.W. Spong, S. Hutchinson, M. Vidyasagar,
  \textit{Robot Modeling and Control},
  Wiley, 2006.

\bibitem{hogan1985}
  N. Hogan,
  ``Impedance Control: An Approach to Manipulation,''
  \textit{ASME J. Dyn. Sys., Meas., Control},
  vol. 107, 1985.

\bibitem{colgate1993}
  J.E. Colgate, J.M. Brown,
  ``Factors affecting the Z-width of a haptic display,''
  in \textit{Proc. IEEE ICRA}, 1994, pp. 3205--3210.

\bibitem{murray1994}
  R.M. Murray, Z. Li, S.S. Sastry,
  \textit{A Mathematical Introduction to Robotic Manipulation},
  CRC Press, 1994.

\bibitem{niemeyer2008}
  G. Niemeyer, C. Preusche, G. Hirzinger,
  ``Telerobotics,''
  in \textit{Springer Handbook of Robotics},
  Springer, 2008.

\bibitem{thompson2001}
  S. Thompson, I. Reid, \textbf{A. Muñoz}, D.W. Murray,
  ``Using a stationary camera to improve the performance
  of a teleoperated robot,''
  \textit{IEEE Trans. Systems, Man \& Cybernetics-A},
  vol. 31(1), pp. 74--82, 2001.

\bibitem{ref9}
  %% COMPLETAR: Citar un artículo IEEE de los últimos 5 años sobre
  %% teleoperación con sensado de fuerza.
  %% Buscar en IEEE Xplore: "force-sensing teleoperation robot arm 2020-2025"
  A.~Autor, B.~Autor, ``Título del artículo,''
  \textit{IEEE Trans. Robotics}, vol.~XX, no.~Y, pp.~ZZ--ZZ, 202X.

\bibitem{ref10}
  %% COMPLETAR: Citar un artículo IEEE reciente sobre estimación
  %% de fuerza sin sensor externo (torque-based force estimation).
  C.~Autor, D.~Autor, ``Título del artículo,''
  in \textit{Proc. IEEE ICRA}, 202X, pp.~ZZ--ZZ.

\end{thebibliography}

%% ============================================================
%% APÉNDICE — Rúbrica de Evaluación Completa
%% ============================================================
\appendix
\section{Rúbrica de Evaluación}
\label{app:rubrica}

\begin{center}
\small
\begin{tabular}{p{0.35\columnwidth}ccp{0.25\columnwidth}}
  \toprule
  \textbf{Pregunta / Criterio} & \textbf{Pts.} & \textbf{Mín.} & \textbf{Evidencia Requerida} \\
  \midrule
  \multicolumn{4}{l}{\textbf{Teoría y Fundamentos}} \\
  T1: Jacobiano y singularidades xArm & 8 & 5 & Derivación + figura sing. \\
  T2: Control de par calculado (CTC)  & 10 & 6 & Derivación + gráfica resp. escalón \\
  T3: Control de impedancia           & 8 & 5 & Cálculo $K_{d,max}$ + tabla $K_d$ \\
  \midrule
  \multicolumn{4}{l}{\textbf{Arquitectura y Comunicación}} \\
  A1: Diagrama de bloques completo    & 10 & 6 & Figura con señales dimensionadas \\
  A2: Protocolo y métricas de red     & 8 & 5 & Tabla RTT/jitter + captura ping \\
  \midrule
  \multicolumn{4}{l}{\textbf{Sensado de Fuerza}} \\
  F1: Estimador torque-fuerza         & 12 & 7 & Gráficas antes/después filtro \\
  F2: Comparativa estrategias         & 8 & 5 & Tabla completada + justificación \\
  \midrule
  \multicolumn{4}{l}{\textbf{Control e Implementación}} \\
  C1: Control del esclavo             & 10 & 6 & Código funcional + flowchart \\
  C2: Feedback háptico en maestro     & 8 & 5 & Cálculo $\beta_{max}$ + código \\
  I1: Setup y configuración xArm SDK  & 10 & 6 & Documentación reproducible \\
  \midrule
  \multicolumn{4}{l}{\textbf{Resultados Experimentales}} \\
  E1: Setup y descripción experimentos & 6 & 4 & Tabla hardware + descripción \\
  E2: Error de seguimiento (6 joints)  & 12 & 7 & Gráficas + tabla métricas \\
  E3: Estimación y validación fuerza   & 14 & 9 & Gráficas + validación báscula \\
  E4: Análisis estabilidad háptica     & 8 & 5 & Prueba hard-wall + $\beta_{crit}$ \\
  \midrule
  \multicolumn{4}{l}{\textbf{Análisis y Conclusiones}} \\
  D1: Discusión limitaciones y EA     & 8 & 5 & Citas IEEE + comparativa cuant. \\
  Cn1: Conclusiones                   & 6 & 4 & Párrafos 200--300 palabras \\
  \midrule
  \multicolumn{4}{l}{\textbf{Calidad del Documento IEEE}} \\
  Formato y ortografía                & 4 & 2 & LaTeX compilable, sin errores \\
  Figuras con pie de figura y etiquetas & 4 & 2 & Mín. 8 figuras propias \\
  Referencias (mín. 10, 2 recientes)  & 4 & 2 & BibTeX correcto \\
  \midrule
  \rowcolor{tec_blue!15}
  \textbf{TOTAL} & \textbf{160} & \textbf{96} & \\
  \bottomrule
\end{tabular}
\end{center}

\begin{rubricbox}[Criterio de aprobación]
  \begin{itemize}[itemsep=0.2em]
    \item Score $\geq 96$ puntos sobre 160 (60\%).
    \item La sección E3 debe incluir \textbf{evidencia experimental real}
          con los xArm físicos (no simulación).
    \item El estimador de fuerza debe mostrar coherencia con los torques
          articulares medidos: $\bm{\tau}_{meas} \approx \bm{J}^\top \hat{\bm{F}}_{ext}$.
    \item El documento LaTeX debe compilar sin errores.
    \item Mínimo \textbf{8 figuras originales} (no tomadas de internet).
  \end{itemize}
\end{rubricbox}

\section{Guía de Entrega}
\label{app:entrega}

\begin{todobox}[Checklist de entrega]
  Antes de entregar, verifiquen que el repositorio/zip contiene:
  \begin{enumerate}[label=$\square$, itemsep=0.3em]
    \item \texttt{reporte.pdf} — artículo IEEE compilado (este documento completo)
    \item \texttt{reporte.tex} — fuente LaTeX del artículo
    \item \texttt{master\_control.py} — código del robot Maestro
    \item \texttt{slave\_control.py} — código del robot Esclavo
    \item \texttt{force\_estimator.py} — módulo de estimación de fuerza
    \item \texttt{net\_test.py} — herramienta de test de red (adaptado para xArm)
    \item \texttt{data/} — carpeta con los archivos \texttt{.csv} o \texttt{.npy}
          de cada ensayo experimental (mínimo 3 ensayos)
    \item \texttt{figures/} — imágenes originales (PNG/PDF) usadas en el reporte
    \item \texttt{video\_demo.mp4} — video de 60--120 s mostrando el sistema
          funcionando en tiempo real con ambos xArm
    \item \texttt{README.md} — instrucciones para reproducir los experimentos
  \end{enumerate}
\end{todobox}

\begin{rubricbox}[Video de demostración — Puntos extra]
  Un video de demostración de alta calidad que muestre:
  \begin{itemize}[itemsep=0.2em]
    \item El operador moviendo el Maestro y el Esclavo replicando el movimiento.
    \item Un evento de contacto claro con las gráficas de fuerza en pantalla.
    \item La reacción háptica perceptible en el Maestro.
  \end{itemize}
  Un video bien editado (subtítulos, señales anotadas) puede obtener
  hasta \textbf{+10 puntos extra} a criterio del profesor.
\end{rubricbox}

\end{document}
%% ============================================================
%% FIN DEL DOCUMENTO
%% ============================================================
